\documentclass[11pt,a4paper]{article}
\usepackage[a4paper,margin=25mm]{geometry}
\usepackage[T1]{fontenc}
\usepackage{lmodern}
\usepackage{amsmath,amssymb,amsthm,mathtools}
\usepackage{booktabs}
\usepackage{microtype}
\usepackage[colorlinks=true,linkcolor=blue,urlcolor=blue,citecolor=blue]{hyperref}
\hypersetup{pdftitle={A cookbook of irrationality proofs, and where to push next},
  pdfauthor={Mathathon peer-review digest}}

\newcommand{\Z}{\mathbb Z}
\newcommand{\Q}{\mathbb Q}
\newcommand{\R}{\mathbb R}
\newcommand{\C}{\mathbb C}
\newcommand{\dd}{\,d}
\newcommand{\Li}{\operatorname{Li}}
\newcommand{\cA}{\mathcal A}
\newcommand{\cB}{\mathcal B}
\newcommand{\st}[1]{\textsf{[#1]}}
\DeclareMathOperator{\lcm}{lcm}
\newtheorem{theorem}{Theorem}
\newtheorem{lemma}[theorem]{Lemma}
\newtheorem{proposition}[theorem]{Proposition}
\newtheorem{corollary}[theorem]{Corollary}
\theoremstyle{definition}
\newtheorem*{problem}{First problem}
\setlength{\parskip}{0pt plus 1pt}
\frenchspacing
\makeatletter
\renewcommand\section{\@startsection{section}{1}{\z@}{-2.2ex plus -.6ex minus -.2ex}{1.1ex plus .2ex}{\normalfont\large\bfseries}}
\renewcommand\subsection{\@startsection{subsection}{2}{\z@}{-1.6ex plus -.5ex minus -.2ex}{.7ex plus .2ex}{\normalfont\normalsize\bfseries}}
\renewcommand\paragraph{\@startsection{paragraph}{4}{\z@}{1.2ex plus .4ex minus .2ex}{-.8em}{\normalfont\normalsize\bfseries}}
\makeatother
\AtBeginDocument{\abovedisplayskip=5pt plus 2pt minus 2pt \belowdisplayskip=5pt plus 2pt minus 2pt
  \abovedisplayshortskip=2pt plus 2pt \belowdisplayshortskip=4pt plus 2pt minus 2pt}

\title{A cookbook of irrationality proofs,\\ and where to push next}
\author{Mathathon peer-review digest\\
{\normalsize (drafted by Claude from the Claude and Codex submissions)}}
\date{14 November 2026}

\begin{document}
{\let\newpage\relax\vspace*{-14mm}\maketitle}\vspace{-3mm}

\noindent\textbf{What this document is.}
Part~I gives one recipe for proving that a number is irrational and runs it, with complete proofs,
on $\log 2$, $e^r$, $\pi^2$ and $\zeta(2)$; the proofs for $\zeta(3)$ are in the companion digest
\path{zeta3-proofs.pdf} (``the companion''). Part~II records why the obvious extensions to
$\zeta(5)$, $\zeta(7)$, Catalan's constant and Euler's constant fail, separating what is proved from
what is only computed. Part~III lists directions to develop, each with a first concrete problem.

\smallskip\noindent\textbf{Status labels.} Every statement carries one of:
\st{P} proved in this document;
\st{C} classical (the references are listed in the last section; none was retrieved);
\st{R} reported by one of the submissions, not verified by us;
\st{N} verified numerically or by exact finite computation only;
\st{Q} question.
All checking so far, ours included, was done by AI agents. No human has refereed anything here, and
we make no claim that any statement is new.

\smallskip\noindent\textbf{Notation.} $D_n=\lcm(1,\dots,n)$, $D_0=1$.
$H_n^{(s)}=\sum_{m=1}^n m^{-s}$ and $H_n=H_n^{(1)}$.
$P_n$ is the shifted Legendre polynomial \eqref{eq:legendre}.
The Ap\'ery numbers are $a_n=\sum_{k}\binom nk^2\binom{n+k}k^2$ $(1,5,73,1445,\dots)$ and $b_n$
$(0,6,\frac{351}4,\frac{62531}{36},\dots)$, the solution of Ap\'ery's recurrence (M4) with $b_0=0$, $b_1=6$;
equivalently $I_n=2(a_n\zeta(3)-b_n)$ for the integral $I_n$ of Section~2.6 (companion, Section~3). As in the companion,
$r_n(t)=\prod_{j=1}^n(t-j)\big/\prod_{j=0}^n(t+j)$, $R_n=r_n^2$, and
$\Lambda=17+12\sqrt2=(1+\sqrt2)^4=33.970563$, so $\Lambda^{-1}=17-12\sqrt2=0.0294373$.
Dictionary: Codex writes $L_n$ for our $P_n$, $A_n,B_n,d_n$ for our $a_n,b_n,D_n$, and
$P_n(z),Q_n(z)$ for the Pad\'e polynomials that we call $U_n(z),V_n(z)$; Codex package~01 writes $P_n(t)$ for our $r_n(t)$. In Fauzan's preprint $D_m(t)$
denotes $\prod_{j\le m}(t+j^2)$, which we write $\Pi_m(t)$.

\section*{\Large Part I. The cookbook}
\section{The recipe}

An irrationality proof of the kind treated here would, if the number were rational, manufacture a positive integer smaller than $1$.

\begin{lemma}[Integer squeeze \st{P}]\label{lem:squeeze}
Let $\theta\in\R$. If there are integers $u_n,v_n$ with $0<|u_n\theta-v_n|\to0$ along an infinite
set of indices $n$, then $\theta$ is irrational.
\end{lemma}
\begin{proof}
If $\theta=p/q$ with integers $p$, $q>0$, then $q(u_n\theta-v_n)=u_np-v_nq$ is a non-zero integer
whose absolute value tends to $0$, which is impossible.
\end{proof}

In practice one constructs rational linear forms $\ell_n=x_n\theta-y_n$ and a positive integer $d_n$
that clears denominators, $d_nx_n,d_ny_n\in\Z$. The lemma then asks for three things.
\begin{description}\itemsep0pt
\item[(I) Integrality with a budget.] A bound $d_n\le e^{(\delta+o(1))n}$.
\item[(N) Non-vanishing.] $\ell_n\ne0$.
\item[(S) Smallness.] $|\ell_n|\le e^{-(\rho+o(1))n}$.
\end{description}
The proof closes when the \emph{balance inequality} $\rho>\delta$ holds. Three remarks guide all that
follows. First, only $d_n\ell_n$ matters, so a good rational approximation $y_n/x_n$ proves nothing by
itself. Second, the three obligations are usually visible in three different representations of
$\ell_n$ (an expansion shows denominators, an integral the sign, an extremal problem or a recurrence
the size), and the work is in passing between them. Third, the lemma needs only infinitely many $n$: if (N) holds for
\emph{every} $n$, the budget (I) may be met along a subsequence. Section~\ref{sec:lcm} exploits this.

\begin{lemma}[Polynomial variant \st{P}]\label{lem:poly}
Let $\theta\in\R$. Suppose there are $Q_n\in\Z[X]$ with $\deg Q_n\le Cn$ and
$0<|Q_n(\theta)|\le e^{-cn^2}$ for infinitely many $n$, where $C,c>0$ are constants. Then
$\theta$ is irrational.
\end{lemma}
\begin{proof}
If $\theta=p/q$, $q>0$, then $q^{\deg Q_n}Q_n(p/q)$ is a non-zero integer of absolute value at most
$q^{Cn}e^{-cn^2}$, which tends to $0$.
\end{proof}
No bound on the coefficients of $Q_n$ is needed. The price of clearing the denominator of $\theta$
is linear in $n$, so it is the \emph{quadratic} scale of the decay that pays. This is the form of the
criterion used by the determinant method of Section~\ref{sec:hankel}.

\section{Worked examples}

\subsection{The tool: shifted Legendre polynomials}
All integral examples need a polynomial with integer coefficients whose integral against a smooth
function is small. Zeros of high order at both endpoints do this through integration by parts, so we
differentiate $x^n(1-x)^n$ exactly $n$ times:
\begin{equation}\label{eq:legendre}
P_n(x)=\frac1{n!}\frac{d^n}{dx^n}\bigl[x^n(1-x)^n\bigr]
=\sum_{k=0}^n(-1)^k\binom nk\binom{n+k}k x^k .
\end{equation}
The expansion follows from $x^n(1-x)^n=\sum_k(-1)^k\binom nkx^{n+k}$ and
$\frac{(n+k)!}{k!\,n!}=\binom{n+k}k$. Thus $P_n\in\Z[x]$. The division by $n!$ costs nothing
arithmetically and is what will make the integrals small.

\begin{lemma}[Rodrigues transfer \st{P}]\label{lem:rod}
If $f$ has $n$ continuous derivatives on $[0,1]$, then
\[
\int_0^1P_n(x)f(x)\dd x=\frac{(-1)^n}{n!}\int_0^1x^n(1-x)^nf^{(n)}(x)\dd x .
\]
In particular $\int_0^1P_n(x)q(x)\dd x=0$ for every polynomial $q$ of degree less than $n$, and
$\int_0^1P_n^2\dd x=\frac1{2n+1}$.
\end{lemma}
\begin{proof}
Integrate by parts $n$ times. Each boundary term contains a derivative of $x^n(1-x)^n$ of order
less than $n$, which vanishes at $0$ and at $1$. For the norm take $f=P_n$: by \eqref{eq:legendre},
$P_n^{(n)}=(-1)^n\binom{2n}n n!$, and $\int_0^1x^n(1-x)^n\dd x=\frac{n!^2}{(2n+1)!}$.
\end{proof}

\subsection{$\log 2$, and $\log\frac{m+1}{m}$ for every integer $m\ge1$}
\emph{Motivation.} We have $\log\frac{m+1}m=\int_0^1\frac{dx}{m+x}$. Integrating a polynomial
$p\in\Z[x]$ against $\frac1{m+x}$ gives a linear form in $1$ and the logarithm, because
$p(x)=p(-m)+(x+m)\,s(x)$ with $s\in\Z[x]$. The coefficient of the logarithm is the integer
$p(-m)$, and the rational part $\int_0^1s$ has denominator dividing $D_{\deg p}$. Lemma~\ref{lem:rod}
says that $p=P_n$ makes the integral small. Set
\[
\mathcal L_n=\int_0^1\frac{P_n(x)}{m+x}\dd x .
\]
\begin{theorem}[\st{C}, proved here]\label{thm:log}
For every integer $m\ge1$ the number $\log\frac{m+1}m$ is irrational.
\end{theorem}
\begin{proof}
\emph{Integrality.} Write $P_n(x)=P_n(-m)+(x+m)s_n(x)$. Since $x^k-(-m)^k$ is divisible by $x+m$
in $\Z[x]$, the polynomial $s_n$ has integer coefficients and degree $n-1$. Hence
$\mathcal L_n=P_n(-m)\log\frac{m+1}m+\rho_n$ with $P_n(-m)\in\Z$ and $D_n\rho_n=D_n\int_0^1s_n\in\Z$.

\emph{Non-vanishing and smallness.} Lemma~\ref{lem:rod} with $f(x)=(m+x)^{-1}$,
$f^{(n)}(x)=(-1)^nn!\,(m+x)^{-n-1}$, gives
\[
\mathcal L_n=\int_0^1\Bigl(\frac{x(1-x)}{m+x}\Bigr)^{\!n}\frac{dx}{m+x}>0 .
\]
The function $h(x)=\frac{x(1-x)}{m+x}$ has $h'(x)=0$ exactly when $(1-2x)(m+x)=x(1-x)$, that is
$x^2+2mx-m=0$, so $x_0=\sqrt{m^2+m}-m$; at that point $h(x_0)=1-2x_0=(\sqrt{m+1}-\sqrt m)^2=:\kappa_m$.
Therefore $0<\mathcal L_n\le\kappa_m^n\log\frac{m+1}m$, and $\kappa_m\le\kappa_1=3-2\sqrt2=0.171573$.

\emph{Balance.} If $\log\frac{m+1}m=p/q$, then $qD_n\mathcal L_n$ is a positive integer, at most
$p\,D_n\kappa_1^n$. By Lemma~\ref{lem:lcmsub} below, $D_n<3^n$ for infinitely many $n$, and
$3\kappa_1=0.514719<1$. Contradiction.
\end{proof}
In the notation of the recipe: $\delta=1$ by the prime number theorem (and $\delta\le\log3$ along a
subsequence by elementary means), $\rho=-\log\kappa_1=1.762747$.

\subsection{The Pad\'e identity and $e^r$}
\emph{Motivation.} Since $e^x$ is its own derivative, integration by parts never stops:
for a polynomial $f$, $\int_0^zf(x)e^x\dd x=\bigl[e^x(f-f'+f''-\cdots)\bigr]_0^z$. So \emph{every}
polynomial $f$ produces an expression (polynomial)$\cdot e^z-$(polynomial). To make it small we want
$f$ small on $[0,z]$; to make the two polynomials have integer coefficients and low degree we want
most endpoint derivatives of $f$ to vanish and the rest to be integers. The choice
$f(x)=x^n(z-x)^n/n!$ does both.

\begin{proposition}[Pad\'e identity for $e^z$ \st{C}, proved here]\label{prop:pade}
Let
\[
U_n(z)=\sum_{j=0}^n\frac{(2n-j)!}{j!\,(n-j)!}\,z^j,\qquad V_n(z)=U_n(-z).
\]
These have integer coefficients, since $\frac{(2n-j)!}{j!(n-j)!}=\binom nj\frac{(2n-j)!}{n!}$. For
every $z\in\C$,
\begin{equation}\label{eq:pade}
V_n(z)e^z-U_n(z)=\tfrac{(-1)^n}{n!}\!\int_0^z\!x^n(z-x)^ne^x\dd x
=\tfrac{(-1)^nz^{2n+1}}{n!}\!\int_0^1\!u^n(1-u)^ne^{zu}\dd u
=z^{n+1}\!\!\int_0^1\!P_n(u)e^{zu}\dd u .
\end{equation}
\end{proposition}
\begin{proof}
Let $z$ be real, $f(x)=x^n(z-x)^n/n!$ and $F=\sum_{j=0}^{2n}(-1)^jf^{(j)}$. Since $f^{(2n+1)}=0$,
$F+F'=f$, so $(e^xF)'=e^xf$ and $\int_0^ze^xf\dd x=e^zF(z)-F(0)$. From
$f(x)=\frac1{n!}\sum_k(-1)^k\binom nkz^{n-k}x^{n+k}$ the only non-zero derivatives at $0$ are
$f^{(n+k)}(0)=(-1)^k\binom nk\frac{(n+k)!}{n!}z^{n-k}$, so, writing $j=n-k$,
$F(0)=(-1)^n\sum_j\binom nj\frac{(2n-j)!}{n!}z^j=(-1)^nU_n(z)$.
The symmetry $f(z-x)=f(x)$ gives $f^{(i)}(z)=(-1)^if^{(i)}(0)$, hence
$F(z)=\sum_if^{(i)}(0)=(-1)^nU_n(-z)$. This proves the first equality for real $z$; the second is
the substitution $x=zu$. For complex $z$ read $\int_0^z$ along the segment, that is, as the second form. The left side
$V_n(z)e^z-U_n(z)$ and the second form are entire functions of $z$ (expand $e^{zu}$ and integrate term
by term) that agree for real $z$, hence for all $z$. The third equality is Lemma~\ref{lem:rod} with
$f(u)=e^{zu}$ (applied to real and imaginary parts).
\end{proof}
The middle expression vanishes to order $2n+1$ at $z=0$ although $U_n,V_n$ have degree $n$: this is
the defining property of the diagonal Pad\'e approximant $U_n/V_n$ of $e^z$. The last expression says
that the Pad\'e remainder is a Legendre moment of the exponential.

\begin{theorem}[\st{C}, proved here]\label{thm:er}
For every non-zero rational $r$, the number $e^r$ is irrational.
\end{theorem}
\begin{proof}
Replacing $r$ by $-r$ replaces $e^r$ by its reciprocal, so let $r=k/m$ with positive integers
$k,m$. \emph{Integrality:} $m^nU_n(r)$ and $m^nV_n(r)$ are integers. \emph{Non-vanishing and
smallness:} for $z=r>0$ the integrand in \eqref{eq:pade} is positive and $x(r-x)\le r^2/4$, so
$0<|V_n(r)e^r-U_n(r)|\le re^r(r^2/4)^n/n!$. \emph{Balance:} if $e^r=p/q$, then
$qm^n\bigl(V_n(r)e^r-U_n(r)\bigr)$ is a non-zero integer of absolute value at most
$q\,re^r(mr^2/4)^n/n!\to0$.
\end{proof}
Here the budget is $m^n$, exponential, and the decay is factorial: the margin is infinite on the
exponential scale, which is why this proof is easy and Ap\'ery's is not.
Consequently $\log q$ is irrational for every positive rational $q\ne1$; Theorem~\ref{thm:log} is
kept because it is the model for $\zeta(2)$ and $\zeta(3)$.

\subsection{$\pi^2$ and $\pi$}
\emph{Motivation.} $e^{i\pi}=-1$ is rational, so the Pad\'e identity at $z=i\pi$ must fail to be an
integer squeeze only because $z$ is not rational. The even part of the identity involves only
$z^2=-\pi^2$. Define
\[
T_n(X)=2\sum_{0\le j\le n/2}\frac{(2n-2j)!}{(2j)!\,(n-2j)!}\,X^j\in\Z[X],
\qquad\text{so that } U_n(z)+V_n(z)=T_n(z^2).
\]
\begin{theorem}[\st{C}, proved here]\label{thm:pi}
$\pi^2$ is irrational; hence so is $\pi$.
\end{theorem}
\begin{proof}
Put $z=i\pi$ in the second form of \eqref{eq:pade}. Since $e^{i\pi}=-1$ and
$(i\pi)^{2n+1}=i(-1)^n\pi^{2n+1}$,
\[
-T_n(-\pi^2)=\frac{i\pi^{2n+1}}{n!}\int_0^1u^n(1-u)^n\bigl(\cos\pi u+i\sin\pi u\bigr)\dd u .
\]
The cosine part vanishes because $u\mapsto1-u$ preserves $u^n(1-u)^n$ and changes the sign of
$\cos\pi u$. Hence
\[
T_n(-\pi^2)=\frac{\pi^{2n+1}}{n!}\int_0^1u^n(1-u)^n\sin(\pi u)\dd u\in\Bigl(0,\ \frac{\pi^{2n+1}}{4^n\,n!}\Bigr].
\]
If $\pi^2=k/m$ with positive integers $k,m$, then $m^{\lfloor n/2\rfloor}T_n(-k/m)$ is a positive
integer, at most $\pi(m^{1/2}\pi^2/4)^n/n!\to0$. Contradiction.
\end{proof}
As far as we recall (Niven 1947, not retrieved), this is Niven's argument with the auxiliary sum of derivatives identified: it is the $e^z$ Pad\'e
approximant at $z=i\pi$. Irrationality of $\pi$ alone would not give that of $\pi^2$.

\subsection{$\zeta(2)$}
\emph{Motivation.} Expanding the geometric series (all terms are non-negative, so sum and integral
may be exchanged),
\begin{equation}\label{eq:z2mom}
\int_0^1\!\!\int_0^1\frac{x^ky^l}{1-xy}\dd x\dd y=\sum_{t\ge1}\frac1{(t+k)(t+l)}=
\begin{cases}\zeta(2)-H_k^{(2)},&k=l,\\[2pt] \dfrac{H_k-H_l}{k-l},&k\ne l.\end{cases}
\end{equation}
So integrating an integer polynomial of degree at most $n$ in each variable against $\frac1{1-xy}$
gives a form in $1,\zeta(2)$ with denominator budget $D_n^2$: the off-diagonal terms have
denominators $D_n\cdot|k-l|$, the diagonal ones $D_n^2$. For smallness use $P_n(x)$ in $x$. The
Rodrigues transfer applied to $x\mapsto\frac1{1-xy}$ produces the factor $y^n$ by itself, since
$\partial_x^n\frac1{1-xy}=\frac{n!\,y^n}{(1-xy)^{n+1}}$; a zero at $y=1$ must be put in by hand. So set
\[
\mathcal Z_n=\int_0^1\!\!\int_0^1\frac{P_n(x)(1-y)^n}{1-xy}\dd x\dd y .
\]
\begin{theorem}[Ap\'ery; proof after Beukers \st{C}, proved here]\label{thm:z2}
$\zeta(2)$ is irrational.
\end{theorem}
\begin{proof}
\emph{Integrality.} Expanding both polynomials and using \eqref{eq:z2mom},
$\mathcal Z_n=\alpha_n\zeta(2)-\beta_n$ with $\alpha_n=\sum_k\binom nk^2\binom{n+k}k\in\Z$ (the diagonal
terms) and $D_n^2\beta_n\in\Z$.

\emph{Non-vanishing and smallness.} For fixed $y\in[0,1)$ apply Lemma~\ref{lem:rod} in $x$, then
integrate over $y$ (the integrand $|P_n(x)|(1-y)^n/(1-xy)$ is integrable on the square, so Fubini
applies):
\[
(-1)^n\mathcal Z_n=\int_0^1\!\!\int_0^1\Bigl(\frac{x(1-x)y(1-y)}{1-xy}\Bigr)^{\!n}\frac{dx\dd y}{1-xy}>0 .
\]
To bound $h=\frac{x(1-x)y(1-y)}{1-xy}$ put $s=\sqrt{xy}$. Since $x+y\ge2s$,
$(1-x)(1-y)\le(1-s)^2$, so $h\le\frac{s^2(1-s)}{1+s}$. The derivative of the right side is
$\frac{2s(1-s-s^2)}{(1+s)^2}$, which vanishes at $g=\frac{\sqrt5-1}2$; using $1-g=g^2$ and
$\frac1{1+g}=g$ the maximum is $g^5=\frac{5\sqrt5-11}2=0.0901699$. Hence
$0<|\mathcal Z_n|\le g^{5n}\zeta(2)$.

\emph{Balance.} If $\zeta(2)=p/q$, then $qD_n^2|\mathcal Z_n|$ is a positive integer, at most $p\,D_n^2g^{5n}$.
Along the infinitely many $n$ with $D_n<3^n$ this is at most $p\,(9g^5)^n$, and
$9g^5=0.811529<1$.
\end{proof}
Rates: $\delta=2$ against $\rho=-5\log g=2.406059$.

\subsection{$\zeta(3)$: summary only}
Replace $\frac1{1-xy}$ by the kernel $\frac{-\log(xy)}{1-xy}$, whose moments are
$2(\zeta(3)-H_k^{(3)})$ on the diagonal and $\frac{H_k^{(2)}-H_l^{(2)}}{k-l}$ off it, and use
$P_n$ in both variables. Then \st{C; proved in the companion, Sections 3--6}
\[
I_n=\int_0^1\!\!\int_0^1\!P_n(x)P_n(y)\frac{-\log(xy)}{1-xy}\dd x\dd y=2(a_n\zeta(3)-b_n),
\quad D_n^3b_n\in\Z,\quad 0<I_n\le2\zeta(3)\Lambda^{-n},
\]
and $27\Lambda^{-1}=0.794806<1$ closes the proof along the indices with $D_n<3^n$. The companion
gives three routes to positivity and smallness: Beukers' triple integral (its Section~4), one
telescoping identity for $R_n$ with a Casoratian, the discrete Wronskian of M4 (Section~5, which reaches only the rate $28^{-n}$,
still enough against $27^n$), and an energy inequality (Section~6). Rates: $\delta=3$ against
$\rho=\log\Lambda=3.525494$.

\section{The lcm toolkit}\label{sec:lcm}
The budget in the integral examples is a power of $D_n$. Three levels of information about $D_n$ are
available; the first is all that the examples above need, because their linear forms are non-zero for
every $n$.

\begin{lemma}[\st{P}; Codex exposition \S6; also companion, Section~2]\label{lem:lcmid}
For every integer $N\ge1$, $\;N!=\prod_{k=1}^{N}D_{\lfloor N/k\rfloor}$.
\end{lemma}
\begin{proof}
Fix a prime $p$. The exponent of $p$ in $D_m$ is the number of $a\ge1$ with $p^a\le m$. So its
exponent in the product is
$\sum_{k}\#\{a\ge1:p^a\le\lfloor N/k\rfloor\}=\sum_{a\ge1}\#\{k:k\le N/p^a\}=\sum_{a\ge1}\lfloor N/p^a\rfloor$,
which is the exponent of $p$ in $N!$.
\end{proof}

\begin{lemma}[\st{P}; same source]\label{lem:lcmsub}
For every real $C>e$ there are infinitely many $n$ with $D_n<C^n$. In particular $D_n<3^n$
infinitely often.
\end{lemma}
\begin{proof}
Suppose $D_m\ge C^m$ for all $m\ge M$. Let $N>M$ and $L=\lfloor N/M\rfloor$. For $k\le L$ we have
$\lfloor N/k\rfloor\ge M$, and all factors in Lemma~\ref{lem:lcmid} are at least $1$, so
\[
N\log N\ \ge\ \log N!\ \ge\ \log C\sum_{k=1}^{L}\Bigl\lfloor\frac Nk\Bigr\rfloor\ \ge\ \log C\,(NH_L-L).
\]
Now $H_L\ge\log(L+1)>\log N-\log M$ and $L/N\le1/M$, so
$\log N\ge\log C\,(\log N-\log M-\frac1M)$. As $\log C>1$ and $M$ is fixed, this fails for large
$N$. Finally $e<3$.
\end{proof}

\noindent The lemma suffices for a budget $D_n^k$ against decay $e^{-\rho n}$ whenever $\rho>k$ and
the form is non-zero for every $n$: here $e\kappa_1=0.4664$, $e^2g^5=0.6663$,
$e^3\Lambda^{-1}=0.5913$.

\smallskip\noindent\textbf{Bounds for all $n$} \st{R: proofs are in the source files; we rechecked
the constants only}. A Chebyshev-type argument with a period-6 step function gives
$\log D_n\le c_6n+O(\log^2n)$ with $c_6=\frac12\log2+\frac34\log3=1.170533$
(Codex package 03, \path{main-explanation.tex}, \S5.2; the referee reports re-derive it).
Chebyshev's period-30 function gives $c_{30}=\frac{14\log2+9\log3+5\log5}{25}=1.105550$ (Codex
package 01, \path{zudilin-explained.md}, \S10). Both are below the
thresholds $\frac13\log\Lambda=1.1752$ and $-\frac52\log g=1.2030$ needed for $\zeta(3)$ and
$\zeta(2)$; $c_6$ is below the $\zeta(3)$ threshold by only $0.0046$. Such bounds are needed when the non-vanishing argument works only
for some $n$ (for instance when a Casoratian shows that one of two consecutive forms is non-zero).

\smallskip\noindent\textbf{The prime number theorem} \st{C}: $\log D_n\sim n$. It is needed only when
the margin $\rho-\delta$ is small, as in the determinant proof for $\zeta(5)$
(Section~\ref{sec:hankel}), and even there a subsequence argument may replace it (direction D4).

\section{Which ingredient does what}
Table~\ref{tab:main} lists, for each proof, which ingredient supplies each obligation of the
recipe. Its last two rows are reported by the sources (Section~\ref{sec:hankel}), not verified by us. The energy row uses the notation of Lemma~\ref{lem:energy}.

\begin{table}[t]
\centering\footnotesize
\setlength{\tabcolsep}{3.5pt}
\begin{tabular}{@{}p{18mm}p{25mm}p{26mm}p{27mm}p{25mm}p{26mm}@{}}
\toprule
constant & object & integrality & non-vanishing & smallness & cost vs.\ decay\\
\midrule
$\log2$ \st{P} & $\int_0^1\frac{P_n(x)}{1+x}dx$ & $P_n\in\Z[x]$; $\int x^j=\frac1{j+1}$ gives $D_n$
 & positive integrand after Rodrigues transfer & $\max\frac{x(1-x)}{1+x}=3-2\sqrt2$ & $1$ vs.\ $1.7627$ (per $n$)\\[2pt]
$e^{r}$, $r=\frac km$ \st{P} & $V_n(r)e^r-U_n(r)$ & integer coefficients; $m^n$ & sign of
 $\int_0^r x^n(r-x)^ne^x$ & the $n!$ from Rodrigues & $n\log m$ vs.\ $n\log n$\\[2pt]
$\pi^2$ \st{P} & $T_n(-\pi^2)$, even part of Pad\'e at $i\pi$ & $T_n\in\Z[X]$, degree
 $\lfloor\frac n2\rfloor$ & $\sin\pi u>0$ on $(0,1)$ & the $n!$ & $\frac n2\log m$ vs.\ $n\log n$\\[2pt]
$\zeta(2)$ \st{P} & $\iint\frac{P_n(x)(1-y)^n}{1-xy}$ & moments \eqref{eq:z2mom}; $D_n^2$ & positive
 integrand after transfer in $x$ & $\max h=g^5$ & $2$ vs.\ $2.4061$\\[2pt]
$\zeta(3)$, Beukers \st{C} & triple integral $=I_n$ & log-kernel moments; $D_n^3$ & positive
 integrand & maximum of the integrand, $\Lambda^{-1}$ & $3$ vs.\ $3.5255$\\[2pt]
$\zeta(3)$, Zudilin \st{C} & $-\sum_{t\ge1}R_n'(t)$ & partial fractions of $R_n$; $D_n^3$ &
 Casoratian of the recurrence & $I_n<\frac6{n^2a_{n+1}}$, growth of $a_n$ & $\log27=3.296$ vs.\ $\log28=3.332$\\[2pt]
$\zeta(3)$, energy \st{C} & $\langle P_n,k_n\rangle=I_n$ & same moments & $a_nI_n=\langle Bf_0,f_0\rangle>0$
 & Cauchy--Schwarz, growth of $a_n$ & $3$ vs.\ $3.5255$\\[2pt]
$\zeta(3)$, Hankel \st{R} & $Q_K\in\Z[X]$ of degree $K$ & $p$-adic valuation of a determinant &
 Gram determinant of a positive weight & logarithmic energy, $\frac92\log2$ &
 $\frac94c_6=2.634$ vs.\ $3.119$ (per $K^2$)\\[2pt]
$\zeta(5)$, Fauzan \st{R} & $Q_n\in\Z[X]$ of degree $37n$ & $p$-adic valuation of the whole determinant,
 summed by PNT & Gram determinant of a positive weight & energy bound, computer-assisted &
 $1.3496$ vs.\ $1.3670$ (per $K^2$, $K=40n$)\\
\bottomrule
\end{tabular}
\caption{The three obligations in each proof, and what each step contributes.}\label{tab:main}
\end{table}

\section{The shared mechanisms}\label{sec:mech}

\paragraph{M1. Endpoint zeros give orthogonality, and orthogonality gives Mellin zeros.}
Lemma~\ref{lem:rod} deduced orthogonality of $P_n$ to $1,x,\dots,x^{n-1}$ from the zeros of
$x^n(1-x)^n$ at $0$ and $1$. Applying the same lemma to $f(x)=x^{t-1}$ (integers $t\ge1$, or real $t\ge n+1$) and using
$\int_0^1x^{t-1}(1-x)^n\dd x=\frac{n!}{t(t+1)\cdots(t+n)}$ gives \st{P; checked numerically}
\[
\int_0^1x^{t-1}P_n(x)\dd x=(-1)^n\frac{(t-1)(t-2)\cdots(t-n)}{t(t+1)\cdots(t+n)} .
\]
The left side is $\sum_kc_k/(t+k)$ with $c_k$ the coefficients of $P_n$, so this is an identity of
rational functions. The zeros at $t=1,\dots,n$ \emph{are} the orthogonality relations, and the
poles at $0,-1,\dots,-n$ record the degree. The right side is $(-1)^nr_n(t)$, and its square is
Zudilin's $R_n(t)$. Summing $(xy)^{t-1}$ over $t\ge1$ gives $\frac1{1-xy}$, and $-\frac d{dt}$ inserts
$-\log(xy)$: this is the dictionary between the integral and the rational-function proofs of
Ap\'ery's theorem.

\paragraph{M2. Invert on polynomials, then use a boundary identity.}
On polynomials of degree at most $N$ the operator $\partial=\frac d{dx}$ is nilpotent, so $(1+\partial)^{-1}=\sum_j(-1)^j\partial^j$ is a
finite sum. The function $F$ in the proof of Proposition~\ref{prop:pade} is $(1+\partial)^{-1}f$, and
$(e^xF)'=e^xf$ is the boundary identity. For $\pi^2$ the corresponding operator is $(\pi^2+\partial^2)^{-1}$
and the boundary identity is Green's formula against $\sin\pi x$. When the operator \emph{raises} the degree
instead (as for $\zeta(3)$), the inverse exists up to a one-dimensional obstruction, and that
obstruction is the Ap\'ery number $a_n$; see direction D3.

\paragraph{M3. Three sources of positivity.}
(i) The Rodrigues transfer replaces the sign-changing $P_n$ by the positive $x^n(1-x)^n$ at the price
of differentiating the kernel; it works when the derivatives of the kernel have a fixed sign
($\log$, $\zeta(2)$, Beukers' integral). (ii) A Gram determinant $\det[\int t^{i+j}\dd\nu]_{i,j<h}$
of a positive measure $\nu$ with infinite support is positive, because the matrix is the matrix of
the quadratic form $c\mapsto\int(\sum c_it^i)^2\dd\nu$. It equals
$\frac1{h!}\int\prod_{i<j}(t_i-t_j)^2\prod_i\dd\nu(t_i)$ \st{C: Andr\'eief's identity}, and the
$\binom h2$ factors of the squared Vandermonde are the source of the $e^{-ch^2}$ scale in
Lemma~\ref{lem:poly}. (iii) An energy identity $a_nI_n=\langle Bf_0,f_0\rangle$ with a positive
operator $B$ (direction D3) gives sign and size at once.

\paragraph{M4. Casoratians are discrete Wronskians.}
If $u_n,v_n$ satisfy the same three-term recurrence, their Casoratian $W_n=u_nv_{n+1}-u_{n+1}v_n$
satisfies a first-order recurrence, exactly as the Wronskian of a second-order differential equation
does (Abel's formula). For Ap\'ery's recurrence
$(n+1)^3u_{n+1}=(2n+1)(17n^2+17n+5)u_n-n^3u_{n-1}$ one gets $(n+1)^3W_n=n^3W_{n-1}$, so
$a_nb_{n+1}-a_{n+1}b_n=\frac6{(n+1)^3}$ \st{P given the recurrence for $a_n$ and $b_n$ (companion, Section~5); checked exactly for $n<20$}. This single identity gives
non-vanishing (two consecutive forms cannot both vanish), monotonicity of $b_n/a_n$, and the exact
error $I_n=2a_n\sum_{m\ge n}\frac6{(m+1)^3a_ma_{m+1}}$ \st{N: 60 digits, $n=0,1,3$; proved in the
companion, Section~5}. The companion's bound $I_n<\frac6{n^2a_{n+1}}$ was checked by us for
$n<30$. The Pad\'e polynomials $U_n$ and $V_n$ both obey $u_{n+1}=2(2n+1)u_n+z^2u_{n-1}$ \st{N: symbolic check, $n\le7$; stated in Codex's
\path{shared-coefficients-exp-pi-zeta.md}}; then $W_n=V_nU_{n+1}-V_{n+1}U_n$ satisfies
$W_n=-z^2W_{n-1}$ and $W_0=2z$, so $W_n=2z(-z^2)^n$ \st{P given the recurrence}.

\emph{Where Ap\'ery's recurrence itself comes from} \st{P in the companion, Section~7; the identity
below rechecked symbolically by us for $n\le5$ and by a referee subagent}. Keep two indices,
$r_nr_m$ in place of $r_n^2$. The mixed family satisfies a three-term relation in $n$ with middle
coefficient $(2n+1)(n(n+1)+2m(m+1)+1)$, and the diagonal is closed by one telescoping identity,
\[
r_n(t)^2+r_{n-1}(t)^2+6\,r_n(t)r_{n-1}(t)=g_n(t+1)-g_n(t),\qquad g_n(t)=\frac4{n^2}(t+n)^2r_n(t)^2 .
\]
On the diagonal $m=n$ the middle coefficient is $(2n+1)(3\lambda+1)$, $\lambda=n(n+1)$, and eliminating
the off-diagonal terms gives $6(2n+1)(3\lambda+1)-(2n+1)(\lambda+1)=(2n+1)(17\lambda+5)$. Is there a
two-index family with a three-term relation and a diagonal telescope of this kind for the $\zeta(2)$
recurrence, or a three-index family for
forms in $\zeta(5)$ \st{Q, from the referee report on Codex package~01}?

\section*{\Large Part II. Why the obvious extensions fail}
\section{Higher powers of the logarithm}\label{sec:logk}

The kernel $\frac{-\log(xy)}{1-xy}$ produced $\zeta(3)$, so one tries
$K_k=\frac{(-\log xy)^k}{1-xy}$ and
$J_n^{[k]}=\iint P_nP_nK_k$, that is, $\iint P_n(x)P_n(y)K_k(x,y)\dd x\dd y$, so that $J_n^{[1]}=I_n$.

\emph{The arithmetic is perfect} \st{P for the formula; N: recomputed exactly for $k=1,3,5$,
$n\le40$}. The moments of $K_k$ are $(k+1)!\,(\zeta(k+2)-H_r^{(k+2)})$ on the diagonal and
$k!\,\frac{H_r^{(k+1)}-H_s^{(k+1)}}{r-s}$ off it, by the same computation as \eqref{eq:z2mom}. Hence
$D_n^{k+2}J_n^{[k]}\in\Z+\Z\zeta(k+2)$ with coefficient $(k+1)!\,a_n$ of $\zeta(k+2)$: the case
$k=3$ isolates $\zeta(5)$ and $k=5$ isolates $\zeta(7)$, with no other zeta value present.

\emph{The analysis fails.} By M1, $J_n^{[k]}=(-1)^k\sum_{t\ge1}R_n^{(k)}(t)$. The double zeros of
$R_n$ at $t=1,\dots,n$ kill the first derivative there, not the higher ones. Codex's Proposition 16.1
(proof in its Appendix~A; \st{R}, we read the proof but did not check every estimate) states, for
fixed $k\ge2$,
\[
J_n^{[k]}\sim\frac{k(k-1)\,4^{k-2}(\log n)^{k-2}}{n^4},\qquad
J_n^{[3]}=\frac{12\,(H_{n-1}+H_{n+1})}{n^2(n+1)^2}+O\Bigl(\frac{\log n}{n^8}\Bigr).
\]
Our exact evaluation \st{N} gives $J_n^{[3]}=5.80083\cdot10^{-3}$, $4.89328\cdot10^{-4}$,
$3.81758\cdot10^{-5}$ for $n=10,20,40$, against the main term $5.80051\cdot10^{-3}$,
$4.89326\cdot10^{-4}$, $3.81758\cdot10^{-5}$; for comparison $I_n=2.2\cdot10^{-17}$,
$4.1\cdot10^{-33}$, $3.5\cdot10^{-64}$. Since $D_n^5\ge n^5$, the quantity $D_n^5J_n^{[3]}$ tends to
infinity. So this extension fails for a reason proved in the source (Codex, Appendix~A; not fully checked by us) and confirmed numerically for $n\le40$: decay $n^{-4}$ instead of exponential.

\emph{Repairs tried} \st{R, numerical only; Codex \S16.4 and \S17}. Removing the first $M$ terms of the
series replaces $P_n$ by $x^MP_n$ and the budget by $D_{n+M}^{k+2}$; polynomials with fourth-order
Mellin zeros have degree $2n$ and budget $D_{2n}^5$, and lose positivity. In the cases computed with $n\ge2$ the
resulting integer forms are far larger than $1$ (for example $8.8\cdot10^{20}$ at $n=10$ in the
second family). These are finite experiments.

\section{The determinant method}\label{sec:hankel}

All statements in this section are \st{R}: they come from Claude's research notes
(\path{fauzan-status.md}, \path{zeta5-margin.md}, \path{hankel-zeta3.md},
\path{zeta7-status.md}, \path{anand-referee.md}), which describe a preprint and two Lean
projects that we did not rebuild. We recomputed only the arithmetic of the displayed constants.

\paragraph{How the $\zeta(5)$ preprint works.}
A.~Fauzan's preprint ``$\zeta(5)$ is irrational'' (17 Sep 2026, doi:10.5281/zenodo.22826419) uses
Lemma~\ref{lem:poly}: with $K=40n$, $N=3n$, $h=37n$ it constructs $Q_n\in\Z[X]$ of degree $37n$
with $0<Q_n(\zeta(5))<\exp(-139n^2/5)$ for large $n$. The construction starts from a $\Q$-linear
functional $\mu_X$ on rational functions of $t$ with simple poles at $t=-j^2$, with rational values
on polynomials and
$\mu_X\bigl(\frac1{t+j^2}\bigr)=j^4(X-H_j^{(5)})-\frac14+\frac1{2j}$, and takes
$Q_n$ proportional to $\det\bigl[\mu_X\bigl(\Pi_N(t)^6t^{i+j}/\Pi_K(t)\bigr)\bigr]_{0\le i,j<h}$.
\emph{Positivity:} at $X=\zeta(5)$ the functional is integration
against the positive weight $w(y)=\frac{(2\pi)^4}{12}y^5\sum_{\ell\ge1}\ell^4e^{-2\pi\ell y}$ in the
variable $t=y^2$, so the determinant is a Gram determinant (M3(ii)). \emph{Size:} by Andr\'eief's
identity and a weighted logarithmic-energy estimate, $\log\det\le UK^2+O(K\log K)$ with
$U=-\frac{2733991}{2000000}=-1.3670$; this step is computer-assisted. \emph{Arithmetic:} $p$-adic
lower bounds for the valuation of the determinant as a whole (far better than entry by entry),
summed over primes with the prime number theorem, give a normaliser of size $e^{A K^2}$ with
$A=A_{200}\approx1.3496$ (the cost constant for the preprint's cutoff parameter $M=200$). The margin is $A_{200}+U\approx-0.0174$, about $1.3\%$ of either term.
The notes report two Lean formalisations (one with the prime number theorem as its
only axiom; one importing it from a project containing \texttt{sorry}s that are reported, not verified,
to be unreached), which repair at least two gaps in the printed proof; neither was built by the
authors of the notes.

\paragraph{An improved arithmetic bound (``Theorem W'').}
The preprint uses one valuation bound (``inner'') for primes $p\le\frac K3$ and another (``outer'')
for $p>\frac K3$. \path{zeta5-margin.md} reports a proof that the inner bound holds for all primes
$p\ge5$ with $p^2>104n+6$ and beats the outer bound for $\frac K3<p<\frac{69}{170}K$. The cost
constant drops by $\frac{3203}{163200}=0.019626$, so $A_{200}$ becomes $1.329961$ and the margin
$-0.037034$ ($2.13$ times the original margin). Status as reported: proof by a subagent, re-derived by the
orchestrating agent; exact checks at $K=80,120,160,200$ with no violation; constants confirmed by an
independent recomputation that did not review the proof; a uniform $O(1)$ term needed for the
asymptotic constant is not written out.

\paragraph{The same method at $s=3$.}
\path{hankel-zeta3.md} carries out the method for $\zeta(3)$ with no numerator and full size:
$\Delta_K(X)=\det[\mu_X(t^{i+j}/\Pi_K(t))]_{i,j<K}$ with
$\mu_X(\frac1{t+j^2})=j^2(X-H_j^{(3)})-\frac12+\frac1{2j}$. Reported: $\deg\Delta_K=K$ with leading
coefficient $\pm K!^2$; $\Delta_K(\zeta(3))>0$; a normaliser $\Phi_K$ with
$\Phi_K\Delta_K\in\Z[X]$ and $\log\Phi_K\le\frac94\varkappa K^2+o(K^2)$ whenever
$\sum_{p\le u}\log p\le\varkappa u+o(u)$; and
$\log\Delta_K(\zeta(3))=-\frac92\log2\cdot K^2+O(K\log K)$, where $\frac92\log2=3.1192$ is the
energy of an explicit equilibrium measure on $[0,\frac34]$ (the upper bound depends on one subagent lemma that the orchestrator did not re-derive). With $\varkappa=c_6$ the exponent is
$-0.485K^2$, so no prime number theorem is needed. The notes call this a proof draft and a variant of known work
(Zudilin 2017; Fauzan's method at $s=3$).

\paragraph{Why $\zeta(7)$ appears out of reach of this family.}
On the $K^2$ scale the decay does not depend on $s$: the weight for $\zeta(s)$ differs from that for
$\zeta(3)$ by a power of $y$, which is invisible at this scale. The arithmetic cost grows roughly
linearly in $s$, because each node $j$ carries $H_j^{(s)}$ and hence $p^{-s}$ for every prime
$p\le j$. (This is heuristic in the notes: the lower bound for the determinant is proved for all odd
$s$, the matching upper bound only for $s=3$.) For the plain family ($\Delta_K$ as above with $H^{(s)}$) the notes report:
\begin{itemize}\itemsep0pt
\item \emph{Proved} (analytic half re-derived by the orchestrator, arithmetic half not): for every
odd $s\ge15$, $\log|P_K(\zeta(s))|\ge[\frac s2(\frac M{M+1})^2-\frac92\log2-o(1)]K^2$ for the primitive
integer polynomial $P_K$ (not a Legendre polynomial) proportional to $\Delta_K$, where $M=\max\{m:2m^2-m+1<s\}$. The bracket is
$+0.214$ at $s=15$. So no integer multiple of the plain determinant is small at $\zeta(s)$. The proof
uses the prime number theorem.
\item \emph{Open}: $s=7,9,11,13$, where the proved arithmetic constants $0.875$, $2.0$, $2.44$,
$2.89$ fall short of $3.119$. Exact data suggest the true cost is
$\frac{s(s+4)}{2(s+2)}K^2$ ($=4.28K^2$ at $s=7$), which would give the no-go for all $s\ge7$.
\item \emph{Numerical}: about $420$ exact determinants in eight modified families at $s=7$
(numerators, other node sets, mixtures with the $\zeta(5)$ functional, derivative series with multiple poles, minors)
all fail; the best cost-to-decay ratio is $1.218$, in the numerator family.
\end{itemize}
A Zenodo preprint (P.~Anand, record 22920911) claims $0<Q_n(\zeta(7))<e^{-19000n^2}$ with the same
architecture. \path{anand-referee.md} reports that this is impossible: integrality of the leading
coefficient contradicts the claimed size of the normaliser, exact computation for $n\le4$ gives
values at least $e^{1274},\dots,e^{22661}$, and the error is one wrong value of an integral over
primes. We make no statement about whether $\zeta(7)$ is irrational.

\section{Catalan's constant and Euler's constant}

\paragraph{Catalan's constant $G=\sum_{k\ge0}(-1)^k(2k+1)^{-2}$} \st{R: \path{catalan-status.md};
exact finite computations, lemmas not re-derived}. Positive functionals of the same type exist; with
integer nodes the weight $\frac{\pi y}4\frac{\sinh\pi y}{\cosh^2\pi y}$ has rational moments and
pole values in $\Q+\Q G$ with no $\pi^2$. The Hankel determinants are positive at $G$, but
$\log|P_K(G)|/K^2$ is $+0.38,+0.46,+0.49,+0.52$ at $K=40,80,120,160$. The obstruction is
quantitative and has two parts. The kernel decays like $e^{-\pi y}$, half as fast as for $\zeta(3)$,
so the decay is only $2\log3=2.197$ per $K^2$ (numerical). The partial sums of the series for $G$ run up to
$2K$, so primes up to $2K$ enter, and the cost is about $\frac83K^2$ (from an observed valuation
pattern and the prime number theorem). Mixed node sets that would cancel $\pi^2$ need at least $2h$
extra nodes for a matrix of size $h$ and gain nothing analytically (a reported lemma). The notes
mention a claimed proof of $G\notin\Q$ by a different determinant construction (OpenAI, Sep 2026);
neither they nor we verified it.

\paragraph{Euler's constant $\gamma$: Sondow's integral} \st{C for the construction, as presented in
Codex \S18; N: the formula below checked to 25 digits for $n\le3$ with Codex's explicit $L_n$, $A_n$}. The integral
$E_n=\iint\frac{[x(1-x)y(1-y)]^n}{(1-xy)(-\log xy)}\dd x\dd y$ is positive and at most
$16^{-(n-1)}E_1$, and $E_n=\binom{2n}n\gamma+L_n-A_n$ with $A_n$ rational and $L_n$ a rational
combination of $\log(n+1),\dots,\log(2n)$. For example $E_1=2\gamma+2\log2-\frac52=0.0407257$
(checked numerically to 30 digits). Positivity and smallness survive; integrality does not, because
the logarithms cannot be cleared. This is an obstruction in obligation (I), the opposite of
Section~\ref{sec:logk}.

\paragraph{Euler's constant without logarithms} \st{N: recomputed by us for $h\le4$}. One has
$\gamma-\frac12=\int_0^\infty\frac{2y}{(e^{2\pi y}-1)(1+y^2)}\dd y$ (checked to 40 digits), and the
weight has rational moments $\int_0^\infty y^{2e}\frac{2y\dd y}{e^{2\pi y}-1}=\frac{(-1)^eB_{2e+2}}{2e+2}$
with Bernoulli numbers $B_m$. The Gram determinants of $t^{i+j}/(t+1)$, $t=y^2$, are therefore
positive rational linear forms in $1,\gamma$ (there is one pole, so the determinant has degree one
in $\gamma$). After removing common factors they are $2\gamma-1$, $33\gamma-19$,
$150980\gamma-87117$, $1576293905340\gamma-909774886661$, with values $0.154$, $0.0481$, $31.0$,
$8.7\cdot10^7$. The rational approximations improve while the integer forms grow: the denominators outrun the
error. Nothing here bears on the irrationality of $\gamma$.

\section*{\Large Part III. Further directions}
\section*{D1. Differential equations: the Ap\'ery operator}

\emph{Idea.} The proof of $e^r\notin\Q$ rests on the differential equation $y'=y$. The Ap\'ery
numbers also satisfy a differential equation, through their generating function, and $\zeta(3)$
appears in it as the one constant that removes a singularity. This gives a definition of the
approximants $b_n/a_n$ that does not mention integrals or rational functions.

\emph{Key formulas} \st{P}. Let $\vartheta=z\frac d{dz}$, so that $\vartheta z^n=nz^n$, and let
\[
\cA(z)=\sum_{n\ge0}a_nz^n,\qquad \cB(z)=\sum_{n\ge0}b_nz^n,\qquad
L=\vartheta^3-z(2\vartheta+1)(17\vartheta^2+17\vartheta+5)+z^2(\vartheta+1)^3 .
\]
For any power series $u=\sum u_nz^n$, the coefficient of $z^{n+1}$ in $Lu$ is
$(n+1)^3u_{n+1}-(2n+1)(17n^2+17n+5)u_n+n^3u_{n-1}$ (with $u_{-1}=0$), and the constant coefficient
is $0$. The sequence $a_n$ satisfies Ap\'ery's recurrence for all $n\ge0$ (at $n=0$ it reads
$a_1=5a_0$), so $L\cA=0$. The sequence $b_n$ satisfies it for $n\ge1$, and at $n=0$ the expression
equals $b_1-5b_0=6$, so $L\cB=6z$. (Checked exactly on the first eight coefficients.) Since
$\vartheta^3=z^3\frac{d^3}{dz^3}+\cdots$, the coefficient of the third derivative in $L$ is
$z^3(1-34z+z^2)$. The finite singular points are therefore $0$ and the two roots
$\Lambda^{\mp1}=17\mp12\sqrt2$ of $z^2-34z+1$; $\infty$ is also singular. The exponents at $0$ are
$0,0,0$, at $z=\Lambda^{-1}$ they are $0,\frac12,1$, and at $\infty$ they are $1,1,1$ (our symbolic
computation).

\begin{proposition}[\st{P}, given the companion's results on $I_n$]\label{prop:ode}
\textup{(a)} $a_n^{1/n}\to\Lambda$ and $b_n^{1/n}\to\Lambda$, so $\cA$ and $\cB$ have radius of
convergence exactly $\Lambda^{-1}=0.0294$.
\textup{(b)} $\zeta(3)\cA(z)-\cB(z)=\frac12\sum I_nz^n$ has radius of convergence exactly
$\Lambda=33.97$ and satisfies $L\bigl(\zeta(3)\cA-\cB\bigr)=-6z$.
\textup{(c)} $\zeta(3)$ is the only number $c$ for which $c\,\cA-\cB$ converges on a disc of radius
greater than $\Lambda^{-1}$.
\end{proposition}
\begin{proof}
From the companion we use $0<I_n\le2\zeta(3)\Lambda^{-n}$ and $I_n^{1/n}\to\Lambda^{-1}$ (Section~4)
and $I_n<6/(n^2a_{n+1})$ for $n\ge1$ (Section~5).
(a) The first term of the series for $I_n$ in M4 gives $I_n\ge\frac{12}{(n+1)^3a_{n+1}}$, so
$a_{n+1}\ge\frac{12}{(n+1)^3I_n}\ge\frac{6\Lambda^n}{\zeta(3)(n+1)^3}$; and $a_{n+1}<\frac6{n^2I_n}$,
whose $n$-th root tends to $\Lambda$. Since $b_n=\zeta(3)a_n-\frac12I_n$ and $I_n\to0$, also
$b_n/a_n\to\zeta(3)$ and $b_n^{1/n}\to\Lambda$.
(b) The radius is $\lim I_n^{-1/n}=\Lambda$, and the equation follows from $L\cA=0$, $L\cB=6z$.
(c) If $c\,\cA-\cB$ and $c'\cA-\cB$ both converged on such a disc, so would $(c-c')\cA$, forcing
$c=c'$ by (a).
\end{proof}
(Numerically $a_{51}/a_{50}=32.98$.) The exponents $0,\frac12,1$ at $\Lambda^{-1}$ suggest a
singularity of square-root type there; that the $\frac12$-solution actually occurs in $\cA$ is not
proved here. In words, the proposition says:
\emph{$\zeta(3)$ is the unique constant for which $c\,\cA-\cB$ continues analytically to a disc beyond
the first singularity $\Lambda^{-1}$, and the continuation reaches exactly as far as the second one,
$\Lambda$.} The approximants $b_n/a_n$ are
ratios of Taylor coefficients of two solutions, and the irrationality proof is the statement that
the Taylor coefficients of the combination that is regular at $\Lambda^{-1}$ are small compared with
$D_n^{-3}$.

\emph{The two classes of functions, in plain words} \st{definitions recalled from memory, not
retrieved}. Roughly, an $E$-function is a series $\sum c_nz^n/n!$ and a $G$-function a series
$\sum c_nz^n$, with rational $c_n$ of at most geometric growth whose common denominators also grow
at most geometrically, satisfying a linear differential equation with polynomial coefficients.
$e^z$ is of the first kind; $\cA$, $\cB$, $\log(1+z)$ and the polylogarithms are of the second. For
$E$-functions the approximants at a rational point carry an $n!$, which is the infinite margin of
Theorem~\ref{thm:er}. For $G$-functions the margin is exponential against exponential and has to be
won by the constants ($3$ against $3.5255$ here); values at rational points near $0$ can be treated
as in Theorem~\ref{thm:log}, and $\zeta(3)$ enters as a \emph{connection constant} between the
singular points $0$ and $\Lambda^{-1}$.
\st{Literature pointers, not retrieved or verified by us: the theory of $E$- and $G$-functions
(Siegel, Shidlovskii, Andr\'e, Chudnovsky, Katz); $L$ as the symmetric square of a second-order
operator; modular parametrisations of $\cA$ (Beukers); lists of Ap\'ery-like or ``Calabi--Yau''
operators (Almkvist--Zagier and others).}

\begin{problem}
(a) Prove, by a telescoping identity as in the companion, Section~5, that
$\alpha_n=\sum_k\binom nk^2\binom{n+k}k$ $(1,3,19,147)$ and the $\beta_n$ of Theorem~\ref{thm:z2}
$(0,5,\frac{125}4,\dots)$ satisfy $(n+1)^2u_{n+1}=(11n^2+11n+3)u_n+n^2u_{n-1}$ \st{C, not proved
here; N: checked for $n\le7$}. Then, as for $L$, the operator
$L_2=\vartheta^2-z(11\vartheta^2+11\vartheta+3)-z^2(\vartheta+1)^2$ kills $\sum\alpha_nz^n$ and
sends $\sum\beta_nz^n$ to $5z$, with singular points $0$, $g^5$, $-g^{-5}$, $\infty$. For $\log2$ the generating function is algebraic,
$\sum_nP_n(-1)z^n=(1-6z+z^2)^{-1/2}$, so the three constants sit on equations of order $1,2,3$.
(b) Starting only from $L\cB=6z$ and $\cB(0)=0$, compute by variation of constants the unique $c$ for
which $c\cA-\cB$ is regular at $\Lambda^{-1}$, and recognise $\zeta(3)$ without summing a series.
Do (b) for $L_2$ first.
\end{problem}

\section*{D2. Matrix exponentials and path-ordered exponentials}

\subsection*{The Pad\'e identity with a matrix argument}
\emph{Idea.} Identity \eqref{eq:pade} is an identity between power series that converge everywhere,
so it remains true when $z$ is replaced by a square matrix $M$:
\begin{equation}\label{eq:padematrix}
V_n(M)e^M-U_n(M)=\frac{(-1)^n}{n!}M^{2n+1}\int_0^1u^n(1-u)^ne^{uM}\dd u .
\end{equation}
Indeed, if two everywhere-convergent power series have the same coefficients $c_j$, substituting
$M$ gives the same matrix $\sum c_jM^j$; on the left the product of the polynomial $V_n$ and the
series $e^z$ becomes the matrix product because powers of $M$ commute, and on the right the series
for $e^{uM}$ may be integrated term by term. (We also checked \eqref{eq:padematrix} numerically for a
rational $3\times3$ matrix and $n\le7$, to 33 digits.) The integer squeeze now applies to all
entries at once.

\begin{theorem}[\st{P}]\label{thm:matrix}
Let $M$ be a $d\times d$ matrix with rational entries. Then $e^M$ has rational entries if and only if
$M$ is nilpotent.
\end{theorem}
\begin{proof}
If $M^d=0$ then $e^M=\sum_{j<d}M^j/j!$ is rational. Conversely, suppose $M$ is not nilpotent and
$e^M=S/q$ with $S$ an integer matrix and $q$ a positive integer. Choose a positive integer $m$ such
that $mM$ is an integer matrix.

\emph{Integrality.} $U_n$ and $V_n$ have integer coefficients and degree $n$, and
$m^nM^j=m^{n-j}(mM)^j$ is an integer matrix for $j\le n$. So
$\mathcal E_n:=q\,m^n\bigl(V_n(M)e^M-U_n(M)\bigr)=m^nV_n(M)\,S-q\,m^nU_n(M)$ is an integer matrix.

\emph{Smallness.} Let $\|\cdot\|$ be the operator norm for the Euclidean norm on $\C^d$. It is
submultiplicative, $\|e^{uM}\|\le e^{u\|M\|}$, and $u(1-u)\le\frac14$, so by
\eqref{eq:padematrix}
\[
\|\mathcal E_n\|\le q\,m^n\,\frac{\|M\|^{2n+1}e^{\|M\|}}{4^n\,n!}\longrightarrow0 .
\]
Every entry of a matrix is bounded in absolute value by its operator norm. Hence there is $n_0$ such
that the integer matrix $\mathcal E_n$ is zero for all $n\ge n_0$.

\emph{Non-vanishing.} Since $M$ is not nilpotent it has an eigenvalue $\lambda_0\ne0$ in $\C$, with
eigenvector $v\ne0$. Then $f(M)v=f(\lambda_0)v$ for every everywhere-convergent power series $f$, so
$\mathcal E_nv=q\,m^n\bigl(V_n(\lambda_0)e^{\lambda_0}-U_n(\lambda_0)\bigr)v$. By \eqref{eq:pade}, $\mathcal E_n=0$
forces
\[
J_n:=\int_0^1u^n(1-u)^ne^{\lambda_0u}\dd u=0\qquad(n\ge n_0).
\]
We show that this is false for large $n$. Let $\omega_n=\int_0^1u^n(1-u)^n\dd u=\frac{n!^2}{(2n+1)!}$
and regard $u^n(1-u)^n/\omega_n$ as a probability density on $[0,1]$, with expectation written
$\mathbb E$. By symmetry $\mathbb E[u]=\frac12$, and
$\mathbb E[u(1-u)]=\frac{\omega_{n+1}}{\omega_n}=\frac{n+1}{2(2n+3)}$, so the variance is
$\mathbb E[u^2]-\frac14=\frac12-\frac{n+1}{2(2n+3)}-\frac14=\frac1{4(2n+3)}$. For $u\in[0,1]$,
$|e^{\lambda_0u}-e^{\lambda_0/2}|=\bigl|\int_{1/2}^u\lambda_0e^{\lambda_0s}\dd s\bigr|\le
|\lambda_0|e^{|\lambda_0|}\,|u-\tfrac12|$. Therefore, by the Cauchy--Schwarz inequality,
\[
\Bigl|\frac{J_n}{\omega_n}-e^{\lambda_0/2}\Bigr|\le|\lambda_0|e^{|\lambda_0|}\,\mathbb E\,|u-\tfrac12|
\le\frac{|\lambda_0|e^{|\lambda_0|}}{2\sqrt{2n+3}}\longrightarrow0 ,
\]
while $|e^{\lambda_0/2}|=e^{\operatorname{Re}\lambda_0/2}>0$. So $J_n\ne0$ for large $n$, a
contradiction.
\end{proof}

\begin{corollary}[\st{P}]\label{cor:matrix}
\textup{(a)} $e^r\notin\Q$ for rational $r\ne0$ \textup{(}$d=1$\textup{)}.
\textup{(b)} For rational $r\ne0$ the numbers $c(r)=\sum_{j\ge0}\frac{(-r)^j}{(2j)!}$ and
$s(r)=\sum_{j\ge0}\frac{(-r)^j}{(2j+1)!}$ are not both rational.
\textup{(c)} $\pi^2$ is irrational.
\textup{(d)} For rational $\theta\ne0$, $\cos\theta$ and $\sin\theta$ are not both rational; so the
angle of a rational point of the unit circle other than $(1,0)$, for instance $\arctan\frac34$, is
irrational.
\end{corollary}
\begin{proof}
For (b) take $M=\bigl(\begin{smallmatrix}0&-r\\1&0\end{smallmatrix}\bigr)$. Then $M^2=-rI\neq0$, so $M$
is not nilpotent, and summing the exponential series gives $e^M=c(r)I+s(r)M$. For $r>0$,
$c(r)=\cos\sqrt r$ and $s(r)=\frac{\sin\sqrt r}{\sqrt r}$. If $\pi^2$ were rational, $r=\pi^2$ would
give $c=-1$, $s=0$, that is $e^M=-I$; this is (c). Taking $r=\theta^2$ gives (d).
\end{proof}

\emph{Honest assessment.} Theorem~\ref{thm:matrix} is a special case of the Lindemann--Weierstrass
theorem \st{C, not retrieved}: if $e^M$ were rational, its eigenvalue $e^{\lambda_0}$ would be
algebraic with $\lambda_0\neq0$ algebraic; and each of $\cos\theta$, $\sin\theta$ is in fact
transcendental. The point is the proof: one page, uniform in $d$, with the same three obligations
as before, and $e^r$ and $\pi^2$ become two cases of one statement. We did not find this formulation
in the sources we read; we have not searched the literature. (Pad\'e approximants of the matrix
exponential are standard in numerical analysis: our recollection, an unverified pointer.)

\begin{problem}
(a) Make the theorem quantitative: bound $\|e^M-S/q\|$ from below in terms of $q$.
(b) The theorem says only that \emph{some} entry of $e^M$ is irrational. For
$M=\bigl(\begin{smallmatrix}0&-1\\1&0\end{smallmatrix}\bigr)$, identity \eqref{eq:padematrix} gives,
for each $n$, two integer linear forms in $1,\cos1,\sin1$ of size about $\frac{n!}{(2n+1)!}$ with
coefficients of size about $\frac{(2n)!}{n!}$; the product of error and coefficient is
$\approx\frac1{2n+1}$ \st{N, $n\le80$} and tends to $0$. Decide whether the forms for $n$ and $n+1$,
together with the Casoratian of M4, prove that $1,\cos1,\sin1$ are linearly independent over $\Q$.
What remains is to show that two of the available forms are independent of a putative relation
$c_0+c_1\cos1+c_2\sin1=0$; for the two forms of one $n$, with $U_n(i)=p+iq$, the determinant is
$c_0(p^2+q^2)+c_1(p^2-q^2)+2c_2pq$.
\end{problem}

\subsection*{Path-ordered exponentials and polylogarithms}
\emph{Idea.} The function $e^{zM}$ solves $Y'=MY$, $Y(0)=I$. When the coefficient matrix depends on
$z$ and its values do not commute, the solution is a ``path-ordered exponential'', a sum of iterated
integrals. Polylogarithms $\Li_k(z)=\sum_{n\ge1}z^n/n^k$ arise this way, because
$\Li_1'(z)=\frac1{1-z}$ and $\Li_k'(z)=\frac{\Li_{k-1}(z)}z$.

\emph{Key formulas} \st{P; verified symbolically}. Let $E_{ij}$ be the $4\times4$ matrix unit and
$N_0=E_{32}+E_{43}$, $N_1=E_{21}$. Both are nilpotent, and $N_0N_1-N_1N_0=E_{31}\ne0$. The system
\[
Y'(z)=\Bigl(\frac{N_0}z+\frac{N_1}{1-z}\Bigr)Y(z)
\quad\text{has the solution}\quad
Y(z)=\begin{pmatrix}1&0&0&0\\ \Li_1(z)&1&0&0\\ \Li_2(z)&\log z&1&0\\
\Li_3(z)&\frac12\log^2z&\log z&1\end{pmatrix}
\]
on $0<z<1$: in each column the second entry has derivative (first entry)$/(1-z)$, and the third and
fourth have derivative (previous entry)$/z$. One has $Y(z)=H(z)\,z^{N_0}$ with $H$ analytic at $0$,
$H(0)=I$ and $z^{N_0}=I+N_0\log z+\frac12N_0^2\log^2z$, so $Y$ is the solution normalised at $0$. At
$z=1$ only $\Li_1(z)=-\log(1-z)$ diverges, and multiplying on the left by
$(1-z)^{N_1}=I+N_1\log(1-z)$ removes it:
\[
\lim_{z\to1^-}(1-z)^{N_1}Y(z)=
\begin{pmatrix}1&0&0&0\\0&1&0&0\\ \zeta(2)&0&1&0\\ \zeta(3)&0&0&1\end{pmatrix}.
\]
So $\zeta(3)$ is an entry of the regularised transport of this system from $0$ to $1$, exactly as
$e^r$ is the transport of $y'=y$ from $0$ to $r$. If $N_0$ and $N_1$ commuted, the solution would be
$z^{N_0}(1-z)^{-N_1}$ and no zeta values would appear; they come from the commutator.

\emph{Is Ap\'ery's construction the Pad\'e approximation for this system?} \st{Q, with a partial
answer.} For $e^z$ the Pad\'e problem is: find polynomials $U,V$ of degree $n$ with
$Ve^z-U=O(z^{2n+1})$; there are $2n+2$ unknowns and $2n+1$ conditions. Now write
$R_n(t)=r_n(t)^2=\sum_{k=0}^n\bigl[\frac{\mathsf a_k}{(t+k)^2}+\frac{\mathsf b_k}{t+k}\bigr]$ with
$\mathsf a_k=\binom nk^2\binom{n+k}k^2$, and put $p_2(w)=\sum\mathsf a_kw^k$, $p_1(w)=\sum\mathsf b_kw^k$. Using
$\sum_{t\ge1}\frac{z^t}{(t+k)^s}=z^{-k}\bigl(\Li_s(z)-\sum_{m=1}^kz^m/m^s\bigr)$ one finds \st{P}
\begin{align*}
\mathcal S_n(z)&:=\sum_{t\ge1}R_n(t)z^t=p_2(\tfrac1z)\Li_2(z)+p_1(\tfrac1z)\Li_1(z)+p_0(\tfrac1z),\\
\mathcal T_n(z)&:=-\sum_{t\ge1}R_n'(t)z^t=2p_2(\tfrac1z)\Li_3(z)+p_1(\tfrac1z)\Li_2(z)+\tilde p_0(\tfrac1z),
\end{align*}
with polynomials $p_0,\tilde p_0$ of degree less than $n$. Because $R_n$ has double zeros at
$1,\dots,n$, both $\mathcal S_n$ and $\mathcal T_n$ are $O(z^{n+1})$ at $z=0$; because $R_n(t)=O(t^{-2})$,
$p_1(1)=\sum\mathsf b_k=0$; and then $\mathcal T_n(1)=2a_n\zeta(3)+\tilde p_0(1)=I_n$. Multiplying by $z^n$, this is
a \emph{simultaneous} Pad\'e problem for the pairs $(\Li_2,\Li_1)$ and $(\Li_3,\Li_2)$ with shared
polynomial coefficients of degree at most $n$: $4n+4$ unknown coefficients, and $2(2n+1)+1=4n+3$ linear conditions (two
vanishing orders $2n+1$, and $p_1(1)=0$). \st{N}: for $n\le5$ we solved this linear system exactly;
its solution space is one-dimensional and is Ap\'ery's. \st{Literature pointer, not retrieved:
Hermite--Pad\'e approximants to polylogarithms, attributed to Beukers (1981) and Sorokin.} The
remainders also match in form: by M1, $\mathcal T_n(z)=z\iint P_n(x)P_n(y)\frac{-\log(xy)}{1-xyz}\dd x\dd y$
for $|z|\le1$, to be compared with $z^{n+1}\int_0^1P_n(u)e^{zu}\dd u$ in \eqref{eq:pade}.

\begin{problem}
(a) Prove for all $n$ that the $4n+3$ conditions above are independent, so that the Pad\'e problem
characterises $R_n$. (b) Write the corresponding problem for the $6\times6$ system containing
$\Li_1,\dots,\Li_5$ and determine by exact linear algebra, for small $n$, which combinations of
$\zeta(3)$ and $\zeta(5)$ can be isolated at $z=1$ and what vanishing conditions are left over.
Section~\ref{sec:logk} suggests that isolating $\zeta(5)$ costs most of the vanishing order.
(c) \st{Q} Is there a derivation of the smallness of $\mathcal T_n(1)$ directly from the system, parallel to
the proof of Theorem~\ref{thm:matrix}?
\end{problem}

\section*{D3. The energy method}

\emph{Idea} \st{P for the lemma; the $\zeta(3)$ instance is Codex's proof, presented in the
companion}. Positivity and smallness of a linear form can both come from one positive quadratic form.

\begin{lemma}[\st{P}; from Codex's boundary-pairings note]\label{lem:energy}
Let $\langle\cdot,\cdot\rangle$ be a real inner product on a space of functions containing the
polynomials, and let $P$ be a non-zero polynomial of degree $n$ orthogonal to all polynomials of
lower degree. Let $B$ be a linear operator
and suppose that \textup{(i)} $P=Br+A$ for a polynomial $r$ of degree less than $n$ and a constant
$A$; \textup{(ii)} there is a function $k$ with $Bk=1$; \textup{(iii)}
$\langle Bf,f\rangle\ge c\|f\|^2$ with $c>0$ for $f=r+Ak$. Then
\[
0<A\,\langle P,k\rangle\le\frac{\|P\|^2}{c}.
\]
\end{lemma}
\begin{proof}
Put $f_0=r+Ak$, so $Bf_0=P$ and $f_0\ne0$. Orthogonality gives
$\langle P,f_0\rangle=A\langle P,k\rangle$, and $\langle P,f_0\rangle=\langle Bf_0,f_0\rangle\ge c\|f_0\|^2>0$.
By Cauchy--Schwarz $c\|f_0\|^2\le\|P\|\,\|f_0\|$, hence $\|f_0\|\le\|P\|/c$ and
$\langle P,f_0\rangle\le\|P\|^2/c$.
\end{proof}
Hypothesis (i) is automatic when $B$ raises the degree of polynomials by exactly one: then
$B$ maps polynomials of degree less than $n$ onto a complement of the constants.

\emph{Instance 1: $\zeta(3)$} \st{C; companion, Section~6}. On $L^2(0,1)$ take
$B=\mathcal H+n(n+1)$ with $\mathcal Hf=-(x^2(1-x)f')'+xf$, for which
$\mathcal Hx^j=(j+1)^2x^{j+1}-j(j+1)x^j$; $k_n(x)=\int_0^1P_n(y)\frac{-\log(xy)}{1-xy}\dd y$; $c=n(n+1)$; and
the constant is $A=a_n$. With $\|P_n\|^2=\frac1{2n+1}$ the lemma gives
$0<I_n\le\frac1{n(n+1)(2n+1)\,a_n}$. Numerically \st{N} the ratio of the two sides is $0.6171$, $0.7052$, $0.70709$ at $n=1,10,120$.
It tends to $1/\sqrt2$ \st{P}: the ratio $\rho_m=a_{m+1}/a_m$ satisfies
$\rho_m=\frac{(2m+1)(17m^2+17m+5)}{(m+1)^3}-\frac{m^3}{(m+1)^3\rho_{m-1}}$, with coefficients tending to
$34$ and $1$, and $\rho_m\ge28$ for $m\ge9$ (companion, Section~5); since $x\mapsto34-\frac1x$ is a
contraction on $[28,\infty)$ with fixed point $\Lambda$, $\rho_m\to\Lambda$. The series of M4 then gives
$n^3a_nI_n\to12\sum_{j\ge0}\Lambda^{-2j-1}=\frac{12}{\Lambda-\Lambda^{-1}}=\frac1{2\sqrt2}$ (dominated
convergence).

\emph{Instance 2: $\log\frac{m+1}m$} \st{P}. On $L^2(0,1)$ take $B$ to be multiplication by $x+m$ and
$k=\frac1{x+m}$. Then $P_n=(x+m)s_n+P_n(-m)$ as in Theorem~\ref{thm:log}, $c=m$, and
$f_0=P_n/(x+m)$. The lemma gives
$0<P_n(-m)\,\mathcal L_n=\int_0^1\frac{P_n(x)^2}{x+m}\dd x\le\frac1{m(2n+1)}$, a second proof of positivity
and (with the growth of $P_n(-m)$) of smallness. Numerically the ratio tends to $\sqrt{m/(m+1)}$.

\emph{Instance 3: $\zeta(3)$ again, by multiplication} \st{R: Codex package~04, which attributes the
route to Pr\'evost (1996), not retrieved; verified by a referee subagent; the displayed facts
rechecked numerically by us for small indices}. Let $\dd\nu(t)=\pi^2\operatorname{sech}^2(\pi\sqrt t)\tanh(\pi\sqrt t)\dd t$
on $t>0$. Then $\int\frac{d\nu(t)}{t+a^2}=2\zeta(3,a+\frac12)$ for $a>0$ (Hurwitz zeta function, $\zeta(3,x)=\sum_{j\ge0}(j+x)^{-3}$), which equals
$2(\zeta(3)-H_n^{(3)})$ at $a=n+\frac12$. The polynomials
$Q_k(t)=\prod_{r<k}(t+(r+\frac12)^2)/(k!)^2$ satisfy the Hilbert-matrix identity
\[
\nu(Q_jQ_k)=\frac1{j+k+1}=\int_0^1u^{j+k}\dd u ,
\]
so $u^k\mapsto Q_k$ is an isometry from polynomials in $L^2(0,1)$ to polynomials in $L^2(\nu)$. It
carries $P_n$ to $\Phi_n=\sum_k(-1)^k\binom nk\binom{n+k}kQ_k$, which is therefore orthogonal to
lower degrees with $\nu(\Phi_n^2)=\frac1{2n+1}$, and $\Phi_n(-(n+\frac12)^2)=a_n$. Lemma~\ref{lem:energy}
with $B$ = multiplication by $t+(n+\frac12)^2$ and $c=(n+\frac12)^2$ gives
\[
0<I_n=\frac1{a_n}\int\frac{\Phi_n(t)^2}{t+(n+\frac12)^2}\dd\nu(t)\le\frac4{(2n+1)^3a_n},
\]
the same leading order as Instance~1. The referee's question \st{Q}: for which pairs (measure, basis) is the Gram
matrix a Hilbert or Cauchy matrix? Each such pair transports a classical orthogonal family together
with its norms.

\begin{problem}
Find an instance for another constant. The natural target is $\zeta(2)$: find an operator $B$ and a function $k$
with $Bk=1$ and $\langle P_n,k\rangle$ equal to a form $\alpha_n\zeta(2)-\beta_n$, or show that no
second-order operator raising degrees by one does this. Then explain why the
Cauchy--Schwarz step in Instance~1 loses exactly the factor $\sqrt2$.
\end{problem}

\section*{D4. Pushing the determinant method}

\emph{What a proof for $\zeta(7)$ would need} \st{R}. According to the notes: about one extra unit
of $p$-adic valuation per row at every prime $p\le K$, or an analytic gain of about $20\%$ (model-based numerical evidence, not a theorem).

\emph{Why determinants are cheaper than their entries} \st{P}. If $A,B$ are integer $h\times h$
matrices and $A$ has rank at most $r$, then $p^r\det(B+p^{-1}A)\in\Z$: expand by multilinearity in
the columns; a term using $k$ columns of $A$ carries $p^{-k}$ and vanishes when $k>r$. Clearing
denominators entry by entry would cost $p^h$. The valuation bounds of Section~\ref{sec:hankel} are
refined versions of this remark.

\emph{Removing the prime number theorem} \st{R: Codex exposition, Appendix~D; we read the proof and
rechecked the final arithmetic}. The subsequence idea of Lemma~\ref{lem:lcmsub} extends to weighted
sums over primes: for $g$ piecewise $C^1$ with compact support in $(0,\infty)$ and
$T_g(x)=\frac1x\sum_p(\log p)\,g(p/x)$, one has $\liminf_{n\to\infty}T_g(qn)\le\int_0^\infty g$ for
every fixed integer $q\ge1$. The proof uses only $\sum_{p\le x}\frac{\log p}p=\log x+O(1)$. Applied to
the combined prime cost of the $\zeta(5)$ construction with cutoff $M=1000$, and taking the
preprint's local and real estimates as given, it yields infinitely many $K$ with
$0<Q_K(\zeta(5))<e^{-3K^2/100}$ (the bound evaluates to $-0.031925$; we recomputed this number from
the stated rational constants). An earlier check in the Claude notes (\path{SUMMARY.md}, \S7) of a
version of this normaliser found a cost of $1.3786$, above $-U$; we have not reconciled the two. This suffices for irrationality by Lemma~\ref{lem:poly}, not for an
irrationality measure.

\begin{problem}
In increasing order of difficulty. (a) Write the determinant proofs for $\log2$, $\zeta(2)$ and
$\pi$ along the lines of the $s=3$ proof; the notes report numerical totals $-0.71$, $-0.83$,
$-0.38$ per $K^2$, so there is room. (b) Referee Theorem~W by hand and combine it with the
weighted-prime lemma to obtain a statement for $\zeta(5)$ with an explicit margin and no prime
number theorem. (c) Prove the cost formula $\frac{s(s+4)}{2(s+2)}K^2$ for the plain family, which
would settle $s=7,\dots,13$ negatively for that family.
\end{problem}

\section*{D5. Smaller items worth a look}

\emph{Upper approximants for $\zeta(3)$} \st{R: Codex package~05; N for our checks; P where marked}. The fractions $b_n/a_n$
increase to $\zeta(3)$. Codex's Schur-complement study proves (referee report: correct) that
$q_h=\frac{5b_h-b_{h-1}}{5a_h-a_{h-1}}$ lies \emph{above} $\zeta(3)$, and reports
$q_h-\zeta(3)\sim4\pi^3(3-2\sqrt2)\Lambda^{-2h}$. The identity
$q_h-\frac{b_h}{a_h}=\frac6{h^3a_h(5a_h-a_{h-1})}$ is immediate from the Casoratian of M4 \st{P}.
Also $q_h>\zeta(3)$ for $h\ge1$ \st{P}. Since $5I_h-I_{h-1}=-2(5a_h-a_{h-1})(q_h-\zeta(3))$, this
is equivalent to $I_{h-1}>5I_h$. By the series of M4, $I_{h-1}\ge\frac{12}{h^3a_h}$ (its first
term). The recurrence and $a_{m-1}\le a_m$ give
$\frac{a_{m+1}}{a_m}\ge33-\frac{48m^2+72m+28}{(m+1)^3}\ge\frac{29}2$ for $m\ge1$, so consecutive terms
of the series for $I_h$ decrease by a factor at least $(\frac{29}2)^2$, and
$I_h\le\frac{12}{(h+1)^3a_{h+1}}\cdot\frac{841}{837}$. Hence
$I_{h-1}/I_h\ge\frac{837}{841}\cdot\frac{a_{h+1}}{a_h}\ge\frac{837}{841}\cdot\frac{29}2>5$.
Numerically (checked for $h\le120$) the ratio to the stated asymptotic is
$1.056$, $1.028$, $1.014$, $1.009$ at $h=20,40,80,120$. So there are explicit two-sided rational
brackets $\frac{b_h}{a_h}<\zeta(3)<q_h$ of width $\frac6{h^3a_h(5a_h-a_{h-1})}$. The source of
$q_h$ is an extremal problem (verified by a referee subagent, with exact agreement for
$h\le10$): with
$w_3(y)=y^3\frac{d^2}{dy^2}\frac1{e^{2\pi y}-1}>0$,
\[
q_h-\zeta(3)=\min_{p\text{ monic},\ \deg p=h-1}\int_0^\infty\frac{p(y^2)^2\,w_3(y)}{\prod_{j=1}^{2h}(y^2+j^2)}\dd y>0 .
\]
This is not a new irrationality proof (the form is $I_{h-1}-5I_h$ up to a factor); it adds a
description in which the sign of $q_h-\zeta(3)$ and the monotonicity of $q_h$ are immediate. The same
construction for $\zeta(5)$ gives convergent one-sided rational approximants whose denominators are
observed to grow far too fast \st{R, numerical only}. Problems (the referee's): prove that the minimum equals
$\frac{5b_h-b_{h-1}}{5a_h-a_{h-1}}-\zeta(3)$ directly from the series
$5I_h-I_{h-1}=-12\sum_{t\ge1}\frac d{dt}\bigl[\frac{t\,r_h(t)r_{h-1}(t)}{t+h}\bigr]$ \st{N: $h=1,2$};
and explain the denominator growth at $s=5$, where
only the arithmetic fails.

\section*{Status and sources}
\footnotesize

\noindent\textbf{Sources.} Paths are relative to \texttt{Peer Review/} unless a branch is named.
Sections 1--2, M1--M4: the Codex cookbook folder (its framework, shared-coefficients,
boundary-pairings and pi-powers notes); proofs rewritten and rechecked here.
Section 3: Codex exposition PDF \S6; \path{main-explanation.tex} (package 03) and
\path{zudilin-explained.md} (package 01). Sections 6 and 8 ($\gamma$): Codex exposition PDF
\S16--18, \S21, App.~A. Sections 7, 8 (Catalan), D4: branch \texttt{research-claude},
\path{research/claude/} (\path{SUMMARY.md}, \path{fauzan-status.md}, \path{zeta5-margin.md},
\path{hankel-zeta3.md}, \path{zeta7-status.md}, \path{anand-referee.md}, \path{catalan-status.md},
\path{literature-status.md}, \path{next-steps.md}); Codex exposition PDF \S19, App.~D. D1: the
homogeneous equation is in the framework note; the inhomogeneous equation,
Proposition~\ref{prop:ode} and the exponents were worked out here. D2: written here. D3: package 03
and the boundary-pairings note; Instance~2 was worked out here; Instance~3 and the recurrence
paragraph of M4 follow the reports on Codex packages 04 and 01 in \path{Referee-reports/}. D5:
package 05, its status file and its referee report.

\noindent\textbf{What we verified ourselves} (mpmath at 40--400 digits, sympy, exact rational
arithmetic). Every numerical constant in Part~I; identity \eqref{eq:pade} in its three forms (real
and complex $z$, $n\le5$) and at $z=i\pi$; the integrals, linear forms and denominators for
$\log\frac{m+1}m$ and $\zeta(2)$; the Mellin formula of M1; the Casoratians, the bound
$I_n<6/(n^2a_{n+1})$ and the telescoping identity of M4; moments, integrality and values of
$J_n^{[k]}$ ($k=1,3,5$; $n=10,20,40$); $E_1$ and the four forms in $\gamma$; $L\cA=0$, $L\cB=6z$,
the exponents, $0<I_n\le2\zeta(3)\Lambda^{-n}$ for $n<15$, and the $\zeta(2)$ analogue; identity
\eqref{eq:padematrix}, the variance and the concentration estimate; the polylogarithm system, its
factorisation and its limit at $1$; the simultaneous Pad\'e problem for $n\le5$; the three instances
of Lemma~\ref{lem:energy} (the third for small indices only); the statements about $q_h$ for
$h\le120$; and the arithmetic of the constants quoted in Sections~7 and D4. Proofs marked \st{P} were checked line by line by the drafting agent.

\noindent\textbf{What we did not verify.} Everything labelled \st{R}: the preprints on $\zeta(5)$
(A.~Fauzan, \path{doi:10.5281/zenodo.22826419}) and $\zeta(7)$ (P.~Anand, Zenodo record 22920911), the Lean
projects, Theorem~W, the $s=3$ determinant proof, the no-go theorem, everything about Catalan's
constant, the full proof of Codex's Proposition~16.1, the Chebyshev-type bounds. Literature pointers
and the classical references named in the text (Beukers 1979, Niven 1947, Sondow, Pr\'evost 1996,
Zudilin arXiv:math/0202159 and arXiv:1507.05697) were not retrieved by the drafting agent.

\end{document}
